\section{Discussion}
\label{sec:discussion}

\paragraph{What the theory buys a practitioner.}
Three quantities estimable from data --- two spectra, a mismatch spectrum, and noise levels --- determine (i)~whether human data helps at all (sign of Eq.~\eqref{eq:c1}), (ii)~whether there is a robot-data scale beyond which it must be discarded (Theorem~\ref{thm:h}), and (iii)~what it is worth in robot-data units (Corollary~\ref{cor:kex}).
The mismatch-\emph{location} effect (Table~\ref{tab:harmful}) is, to our knowledge, a new and actionable prediction: retargeting errors that live on low-frequency, task-critical coordinates are far more dangerous than equal-magnitude errors on high-frequency coordinates --- which suggests concrete filtering strategies for egocentric pipelines.

\paragraph{Boundaries of the theory.}
The theory has a sharp negative edge, stated fully in the supplementary material.
If the encoder aligns the two domains in feature space ($W_1$ gap $\varepsilon$) while the action marginals differ (mean gap $|\Delta\mu|$), then \emph{any} deterministic $L$-Lipschitz shared head pays mixture risk $R_\lambda \ge \lambda(1-\lambda)(|\Delta\mu| - L\varepsilon)_+^2$, and the mixture-optimal head pays $R_r \ge \lambda^2(|\Delta\mu| - L\varepsilon)_+^2$ on the robot domain --- a floor that \emph{grows} with the mixing ratio, the squared-loss analogue of the bound of~\cite{zhao2019learning}.
Over-strong alignment losses are likewise predicted to collapse useful within-domain structure, the overlapping/collapse regime of~\cite{lei2026simreal}; the first batch observes the predicted signature --- strong MMD alignment reduces closed-loop coverage at intermediate $\lambda$ relative to both light alignment and naive mixing (Sec.~\ref{sec:experiments}, F4).
Finally, with only two embodiment ``environments'', IRM-style full invariance is unidentifiable~\cite{arjovsky2019invariant}; this is why the framework claims only \emph{partial} invariance (appearance-invariant features with a domain-conditioned action head).

\paragraph{Limitations.}
Our theory is stated for a diagonalized ridge estimator over linear labels; the gap to the full ridge solution is 8--28\% on the weighted risk (a full-ridge extension is stated as a conjecture in the supplementary), and the gap to diffusion policies is bridged only at the level of structural predictions, which is all we claim.
The experiments so far cover one task, one policy class, one surrogate ego domain, and a coarse single-seed protocol; the alignment-strength axis (two MMD weights) and a 150-epoch arm are included, but multi-seed significance, the coverage axis, and action-marginal-shift tests are future work.
The exchange rate prices \emph{samples}, while empirical practice prices \emph{hours}; the conversion depends on collection throughput.
Finally, the harmful-region criterion assumes the shared-spectral-basis structure of Assumption~\ref{ass:a1}; heavy basis misalignment across embodiments is outside the current model.

\paragraph{Outlook.}
The immediate next steps are set by the theory itself: a multi-seed significance pass over the interior optimum, the alignment-strength response (why does strong alignment hurt, and where is the transition?), the open-loop-to-closed-loop bridge (compounding-error bounds in the spirit of~\cite{ross2011reduction}), and a real-robot egocentric pairing of the testbed.
More broadly, we hope the exchange-rate formulation turns "how much egocentric data should I mix?" from folklore into a measurable, computable quantity.

\paragraph{Conclusion.}
We gave a closed-form, numerically exact account of cross-embodiment data mixing: a per-coordinate risk theorem, an exact harmful-mixing phase criterion, and an exchange-rate corollary that quantitatively matches empirical substitution ratios.
A controlled dual-domain Push-T testbed confirms the theory's structural predictions with diffusion policies and surfaces an open-loop/closed-loop divergence as a target for the next iteration of the theory.
